\documentclass{article}
\usepackage{ctex, amsmath, amssymb, array, booktabs, geometry,enumitem,tasks}
\geometry{a4paper, margin=1.5cm}
\usepackage{titling}\setlength{\droptitle}{-2cm}

\author{学号：\underline{\hspace{4cm}} 姓名：\underline{\hspace{4cm}} }
\title{套利定价模型(APT)}
\date{2026年9月30日}
\begin{document}
\maketitle

\begin{enumerate}[label=\textbf{\arabic*.}]

% ---------- 4.1 多因子线性模型 ----------
\item \textbf{（因子模型的基本形式）}
设某风险资产的收益率满足单因子模型
\[
R_i = a_i + b_i F_1 + \varepsilon_i
\]
其中 $E(F_1)=0$，$E(\varepsilon_i)=0$，$\mathrm{Var}(F_1)=1$，$\mathrm{Var}(\varepsilon_i)=S_i^2$，且 $\varepsilon_i$ 与 $F_1$ 独立。
\begin{enumerate}[label=(\arabic*)]
    \item 写出 $E(R_i)$ 的表达式；
    \item 证明 $\mathrm{Var}(R_i)=b_i^2+S_i^2$。
\end{enumerate}

\item \textbf{（敏感系数的含义）}
已知三种资产对同一因子的敏感系数分别为 $b_1=0.5$，$b_2=1.5$，$b_3=2.0$。
\begin{enumerate}[label=(\arabic*)]
    \item 若该因子 $F_1$ 意外上涨 $2\%$，在其他条件不变的情况下，三种资产收益率分别变化多少？
    \item 简要说明 $b_i$ 的绝对值越大意味着什么。
\end{enumerate}

\item \textbf{（投资组合的因子敏感系数）}
设市场上有两种资产，其对两因子的敏感系数为：资产 1 的 $b_{11}=0.8$、$b_{12}=1.2$；资产 2 的 $b_{21}=1.6$、$b_{22}=0.4$。某投资者将资金的 $60\%$ 投入资产 1，$40\%$ 投入资产 2。
\begin{enumerate}[label=(\arabic*), topsep=2pt, itemsep=0pt]
    \item 求该投资组合对因子 1 的敏感系数；
    \item 求该投资组合对因子 2 的敏感系数。
\end{enumerate}

\item \textbf{（组合残差方差的分散化）}
设投资组合由 $n$ 种等权资产构成，即 $w_i=1/n$，各资产残差相互独立，且 $\mathrm{Var}(\varepsilon_i)=S_i^2\le S$。
\begin{enumerate}[label=(\arabic*)]
    \item 写出组合残差方差 $\sigma^2(\varepsilon_P)$ 的表达式，并证明 $\sigma^2(\varepsilon_P)\le S/n$；
    \item 当 $n\to\infty$ 时，$\sigma^2(\varepsilon_P)$ 趋于多少？这说明分散化能消除哪一类风险？
\end{enumerate}

\item \textbf{（组合方差的计算）}
已知某双因子模型中，某投资组合满足
\(
\sum_{i} w_i b_{i1}=0.6,\quad \sum_{i} w_i b_{i2}=0.9,
\) 
且组合残差方差 $\sigma^2(\varepsilon_P)=0.02$，两因子方差均为 $1$ 且相互独立。
\begin{enumerate}[label=(\arabic*)]
    \item 根据式(4.1.7)，求组合方差的近似值；
    \item 若忽略残差项，组合方差约等于多少？
\end{enumerate}

% ---------- 4.2 无残差套利定价 ----------
\item \textbf{（套利组合的判定）}
设市场上有三种股票，其期望收益率与单因子敏感系数为
\[
E(R_1)=15\%,\ E(R_2)=21\%,\ E(R_3)=12\%, \quad 
b_1=0.9,\ b_2=3.0,\ b_3=1.8.
\]
考虑组合权重 $w_1=0.1$，$w_2=0.075$，$w_3=-0.175$。
\begin{enumerate}[label=(\arabic*)]
    \item 验证 $w_1+w_2+w_3=0$ 与 $0.9w_1+3.0w_2+1.8w_3=0$；
    \item 计算组合的期望收益率 $E(R_w)$，并判断是否存在套利机会。
\end{enumerate}

\item \textbf{（套利定价线的确定）}
沿用第 6 题的数据，假设无套利条件下
\(
E(R_i)=\lambda_0+\lambda_1 b_i. 
\)
已知 $\lambda_0=8\%$，$\lambda_1=4\%$。
\begin{enumerate}[label=(\arabic*)]
    \item 分别写出无套利时三种股票的期望收益率；
    \item 与原来的 $E(R_1)=15\%$、$E(R_2)=21\%$、$E(R_3)=12\%$ 比较，指出哪些股票价格被高估、哪些被低估；
    \item 说明套利交易会如何推动价格变动，使套利机会消失。
\end{enumerate}

\item \textbf{（风险溢价 $\lambda_k$ 的经济含义）}
设某 $K$ 因子无残差模型满足无套利条件，$E(R_i)=\lambda_0+\sum_{k=1}^{K}\lambda_k b_{ik}$。
现构造一个特殊组合 $\mathbf{w}$，满足
\[
\sum_i w_i b_{ik}=1,\qquad \sum_i w_i b_{ij}=0\ (j\neq k).
\]
\begin{enumerate}[label=(\arabic*)]
    \item 求该组合的期望收益率 $E(R_w)$；
    \item 由此写出 $\lambda_k$ 的表达式；
    \item 若市场存在无风险资产且 $R_f=4\%$，$\lambda_k=6\%$，求该组合的期望收益率。
\end{enumerate}

\item \textbf{（双因子模型下的期望收益率）}
设某资产的期望收益率满足双因子 APT 模型
\[
E(R_i)=R_f+\lambda_1 b_{i1}+\lambda_2 b_{i2},
\]
已知 $R_f=4\%$，两因子的风险溢价分别为 $\lambda_1=5\%$，$\lambda_2=3\%$。
\begin{enumerate}[label=(\arabic*)]
    \item 若 $b_{i1}=0.8$，$b_{i2}=1.2$，求 $E(R_i)$；
    \item 若另一资产的敏感系数为 $b_{j1}=1.5$，$b_{j2}=0.5$，求 $E(R_j)$；
    \item 比较两资产的期望收益率，说明哪个资产承担了更多的因子风险。
\end{enumerate}

\item \textbf{（反向求解无风险利率）}
考虑一个双因子模型，因子 $a$ 和 $b$ 对应的风险溢价分别为 $\lambda_a=4\%$ 和 $\lambda_b=6\%$。股票 $A$ 在两因子上的敏感系数分别为 $b_{A a}=1.2$，$b_{A b}=0.9$，其预期收益率为 $16\%$。假设不存在套利机会。
\begin{enumerate}[label=(\arabic*)]
    \item 写出股票 $A$ 满足的 APT 定价方程；
    \item 求无风险收益率 $R_f$；
    \item 若另一股票 $B$ 的敏感系数为 $b_{B a}=0.6$，$b_{B b}=1.5$，求其无套利条件下的预期收益率。
\end{enumerate}


\item \textbf{（多因子模型与组合方差综合）}
设某双因子模型中，三种资产对两因子的敏感系数及残差方差如下表：

\begin{table}[h]
    \centering
    \begin{tabular}{cccc}
        \toprule
        资产 & $b_{i1}$ & $b_{i2}$ & $\mathrm{Var}(\varepsilon_i)$ \\
        \midrule
        资产 1 & 1.0 & 0.5 & 0.04 \\
        资产 2 & 0.6 & 1.2 & 0.09 \\
        资产 3 & 1.4 & 0.3 & 0.01 \\
        \bottomrule
    \end{tabular}
\end{table}

已知两因子方差均为 $1$ 且相互独立，各资产残差相互独立。某投资者构造等权组合，即 $w_i=1/3$。
\begin{enumerate}[label=(\arabic*)]
    \item 求组合对因子 1 和因子 2 的敏感系数；
    \item 求组合的残差方差 $\sigma^2(\varepsilon_P)$；
    \item 根据式(4.1.7)，求组合方差的近似值；
    \item 若把资产 3 的权重提高到 $1/2$，其余两资产各占 $1/4$，重新计算组合对两因子的敏感系数，并说明组合风险结构发生了什么变化。
\end{enumerate}

\item \textbf{（套利机会的识别与定价线修正）}
设市场上有四种股票，均只受单因子影响，其期望收益率与敏感系数为
\[
E(R_1)=10\%,\ E(R_2)=14\%,\ E(R_3)=18\%,\ E(R_4)=8\%,
\]
\[
b_1=0.5,\ b_2=1.0,\ b_3=1.5,\ b_4=0.8.
\]
已知无风险利率 $R_f=4\%$。
\begin{enumerate}[label=(\arabic*)]
    \item 以 $R_f$ 和某一敏感系数为 $1$ 的组合为基准，写出无套利定价线 $E(R)=\lambda_0+\lambda_1 b$；
    \item 判断上述四只股票中哪些在定价线上方、哪些在下方，是否存在套利机会；
    \item 对存在套利机会的股票，构造一个零成本、零因子暴露且期望收益为正的套利组合（写出权重并验证三个条件）；
    \item 说明套利交易如何使这四只股票的期望收益率回到同一条定价线上。
\end{enumerate}



\end{enumerate}

\end{document}